\chapter{Autodiff Implementation}

Autodiff in \einlang{} is part of the source language. The user writes a primal
program normally, binds the values of interest, and then asks for tangents or
derivatives with \texttt{@}. This design creates implementation pressure: the
compiler must understand derivative requests while high-level tensor structure
is still available, but the backend should not execute raw source-level
\texttt{@} nodes.

\section{Source Forms}

The two core forms are:

\begin{lstlisting}[style=einlang]
let tx = @x;
let dy_dx = @y / @x;
\end{lstlisting}

\texttt{@x} asks for the tangent of a named binding. The quotient form asks for
the derivative or Jacobian of a target binding with respect to a parameter
binding. Tensor quotients may be large, so the implementation can represent
them lazily instead of immediately materializing dense Jacobians.

These two forms are intentionally name-based. The compiler can attach the
request to DefIds for \texttt{x}, \texttt{y}, and the parameter being
differentiated, which makes the derivative a query over the already-resolved
program rather than a transformation over anonymous source text. That choice
also explains why later passes can lower a derivative request into JVP, VJP, or
lazy-Jacobian IR: the request has a stable target, a stable parameter, and
known shape metadata.

\section{Autodiff IR}

The IR contains explicit nodes for autodiff intent:
\texttt{DifferentialIR}, \texttt{JvpIR}, \texttt{VjpIR}, and
\texttt{LazyJacobianIR}. These nodes are allowed in the high-level compiler
phase but must be removed or lowered before final execution. The separate
\pathcode{src/einlang/passes/autodiff_leak_check.py} pass enforces that phase
boundary.

\section{Compile-Time Differentiation}

\pathcode{src/einlang/passes/autodiff/compiletime.py} implements the main
high-level autodiff pass. It clones IR templates, simplifies generated
expressions, tracks dependency queries, and synthesizes derivative expressions
for supported operations. It runs after range, shape, type, extremum
canonicalization, and pre-autodiff pruning, but before Einstein lowering. That
placement is important because derivative generation can still see source-level
indexed definitions and reductions.

The pass also records graph facts: binding maps, function maps, local contexts,
leaf DefIds, and self-recursive DefIds. The current primary path lowers
requests back into ordinary IR; the recorded graph facts remain available for
diagnostics and for the thin runtime autodiff builtin handler.

\subsection{Source Request Rewrite Algorithm}

The first autodiff pass rewrites source-level request syntax while high-level
tensor IR is still visible. Plain \texttt{@x} is expanded immediately. A
quotient \texttt{@y / @x} becomes a lazy-Jacobian IR node when both sides are
named bindings, because the dense Jacobian may be too large to build eagerly.

\begin{lstlisting}[style=pythoncode,caption={Compile-time autodiff request rewriting.}]
autodiff_pass(program, tcx):
    bindings = clone binding_map(program, tcx)
    local_contexts = clone local bindings inside blocks/functions
    targets = collect DifferentialIR and quotient pairs
    if targets is empty:
        record empty autodiff analysis
        return program

    differentiator = Differentiator(bindings, local_contexts, resolver)
    for statement in program.statements:
        statement = rewrite_expr(statement)
    for function in function_ir_map:
        function.body = rewrite_expr(function.body)
    record cloned graph facts for request lowering and diagnostics

rewrite_expr(expr):
    if expr is DifferentialIR(operand):
        return differentiator.standalone(operand, symbolic=False)
    if expr is DifferentialIR(target) / DifferentialIR(wrt)
       and both are identifiers:
        return LazyJacobianIR(target=target, wrt=wrt)
    recursively rewrite children
\end{lstlisting}

The rewrite pass is the phase boundary between source syntax and differentiable
IR. It first clones the binding and local-context maps because derivative
generation may need to re-enter source expressions after the main traversal has
mutated them. It then rewrites plain differentials and quotient requests while
the high-level tensor structure is still present. That timing matters: once
Einstein lowering has erased source-level clauses into loops, it is much harder
to preserve a derivative as a shaped indexed value.

The lazy-Jacobian case is the most important branch. A quotient such as
\texttt{@y/@x} might describe a scalar gradient, a vector Jacobian, or a large
tensor product of target axes and parameter axes. The rewrite pass does not
decide immediately how much of that value to compute. It records the target and
parameter identities, leaving request lowering to choose a row, column, entry,
or full materialization after it sees how the value is consumed. This is the
implementation reason local derivative requests can remain concise in source.

The clone of binding and local-context maps is not an optimization detail; it is
what keeps rewriting stable. During the traversal, request nodes may be replaced
by generated IR, and blocks may gain derivative-local bindings. If derivative
generation looked only at the partially rewritten program, a later request could
observe a mixture of source and generated structure. Keeping a snapshot gives
the differentiator a consistent primal graph for all requests in the pass.

The rewrite also establishes the shape contract for derivatives. Plain
\texttt{@x} can often be expanded immediately because it asks for the tangent of
one value under the current differentiation context. A quotient has a product
shape formed from target axes and parameter axes, so it is represented lazily
until a use site reveals whether the program needs a row, column, entry, or
full value. The analysis point is that \einlang{} treats a derivative request as
a first-class shaped expression while postponing the expensive part of deciding
how much derivative structure to compute.

\textit{Concrete example.} In \texttt{let loss = sum[i]((pred[i] -
target[i]) * (pred[i] - target[i])); let g = @loss / @W;}, the rewrite pass
does not build the full Jacobian immediately. It records a
\texttt{LazyJacobianIR} whose target DefId is \texttt{loss} and whose
with-respect-to DefId is \texttt{W}. If \texttt{loss} is scalar and \texttt{W}
is a matrix, the lazy value has the shape of \texttt{W}; if a later expression
indexes only one element of \texttt{g}, request lowering can compute that
projection instead of materializing the entire derivative value.

\subsection{Differentiation Rules}

The differentiator handles the ordinary expression structure of the language:
literals produce zero tangents, the differentiated variable produces an identity
tangent, addition and subtraction distribute, multiplication uses the product
rule, division uses the quotient rule, and reductions differentiate through
their bodies. Indexed tensor expressions preserve their index structure where
possible, so a derivative of an indexed definition can itself be represented as
an indexed or lazy shaped value instead of being flattened immediately.

The pass is conservative around unsupported constructs. It is better for a
derivative request to remain unsupported than for the compiler to silently
switch to a finite-difference approximation or erase a dependency.

\subsection{Forward Differentiation Algorithm}

The differentiator computes symbolic forward tangents over IR. A seed map says
which binding is differentiated with respect to which input; every other known
leaf receives a zero seed. For top-level bindings, the differentiator memoizes
derivative expressions by DefId so repeated uses of the same dependency do not
duplicate work.

\begin{lstlisting}[style=pythoncode,caption={Symbolic forward-mode differentiation.}]
differentiate_expr(expr, wrt_defid, seed_override=None):
    seed_map = {}
    seed_map[wrt_defid] = seed_override or ones_like(binding(wrt_defid))
    for leaf in graph_leaf_bindings:
        if leaf.defid != wrt_defid:
            seed_map[leaf.defid] = zeros_like(leaf)
    for binding in top_level_bindings:
        if binding does not depend on wrt_defid:
            seed_map[binding.defid] = zeros_like(binding)
    return simplify(diff(expr, seed_map, local_bindings={}, local_diffs={}))

diff(expr):
    case literal:
        return 0
    case identifier:
        if identifier.defid in seed_map:
            return clone(seed_map[identifier.defid])
        return diff(binding(identifier.defid).expr)
    case a + b:
        return diff(a) + diff(b)
    case a - b:
        return diff(a) - diff(b)
    case a * b:
        return a * diff(b) + b * diff(a)
    case a / b:
        return (b * diff(a) - a * diff(b)) / (b ** 2)
    case a ** literal n:
        return n * (a ** (n - 1)) * diff(a)
    case array[index]:
        return diff(array)[index], substituting indices into Einstein bodies
    case if condition then x else y:
        return if condition then diff(x) else diff(y)
    case block:
        emit primal local bindings and matching tangent bindings in order
    case sum[indices](body):
        return sum[indices](diff(body))
    case max[indices](body):
        return select derivative body at argmax of primal body
    case EinsteinIR clauses:
        differentiate each clause value and drop zero clauses
    case function call:
        inline function body with primal arguments, or use custom @fn rule
\end{lstlisting}

Two implementation details make this work for indexed programs. First,
rectangular access into a differentiated \texttt{EinsteinIR} substitutes access
indices into the clause body instead of materializing the whole tangent tensor.
Second, block differentiation emits tangent bindings beside primal bindings, so
self-recursive recurrence definitions can refer to their own tangent history.

The differentiation algorithm is symbolic and structural. It does not execute
the program to build a tape; it walks IR and constructs new IR that represents
the tangent computation. The seed map gives the differentiated parameter an
identity tangent and gives unrelated leaves zero tangents. Memoization by DefId
prevents repeated dependencies from expanding into duplicated derivative
expressions, which is especially important when several indexed clauses read
the same binding.

The indexed cases are where the implementation differs from a scalar textbook
derivative. A tensor access into a differentiated binding can often be lowered
by substituting the access indices into the derivative of the binding's clause.
That preserves sparsity and avoids materializing tangents that the program
never observes. Reductions also keep their loop variables, so differentiating a
sum produces another sum rather than a flattened vector operation. The result
is still ordinary IR that later range, shape, and lowering passes can process.

The block case is where symbolic differentiation meets imperative-looking
source order. A block may introduce local bindings that later expressions use.
The differentiator therefore emits the primal local and the corresponding
tangent local in sequence, preserving the dependency order of the original
program. This lets derivative expressions for later locals refer to earlier
tangent values by DefId, just as primal expressions refer to earlier primal
values. Without this paired local structure, differentiating blocks would either
duplicate subexpressions aggressively or require a separate runtime tape.

The function-call case is intentionally split between inlining and custom
\texttt{@fn} rules. Inlining keeps ordinary source functions transparent to the
differentiator, which is useful for small standard-library helpers. Custom rules
provide a controlled escape hatch for operations whose derivative is simpler or
more stable than their source expansion. Both paths still return IR, not Python
callbacks. That keeps generated derivative code inside the same range, shape,
type, lowering, and validation pipeline as handwritten \einlang{} code.

\textit{Concrete example.} For \texttt{let y = x * x + 3.0 * x; let dy =
@y / @x;}, the seed map gives \texttt{x} tangent one. The structural rules
produce \texttt{x * 1 + x * 1 + 3.0 * 1}, which simplification can present as
\texttt{2.0 * x + 3.0}. For an indexed sum \texttt{let s = sum[i](w[i] *
w[i]);}, the derivative with respect to \texttt{w[j]} stays sparse: the
generated expression preserves the reduction/index relationship rather than
first flattening \texttt{w} into an anonymous vector.

\section{Request Lowering}

\pathcode{src/einlang/passes/autodiff/runtime_requests.py} lowers remaining
autodiff requests into ordinary IR before backend execution. It
contains logic for basis tensors, identity Jacobians, scalar objectives,
cotangent construction, metadata stamping, and plain request lowering. The pass
can lower projected Jacobian uses without forcing a full dense matrix when a row
or column is enough.

For a lazy Jacobian access, the lowerer splits indices into output axes and
with-respect-to axes. If all output axes are fixed, the access becomes a VJP
with a basis cotangent. If all parameter axes are fixed, the access becomes a
JVP with a basis tangent. If neither side is fixed enough, the lowerer falls
back to materializing the lazy value through a generated indexed expression.

\subsection{JVP and VJP Lowering Algorithm}

Request lowering is still compile-time IR generation. It builds seed tensors,
invokes the same differentiator, simplifies the generated tensor IR, and stamps
type and shape metadata back onto the result.

\begin{lstlisting}[style=pythoncode,caption={Plain IR lowering for JVP and VJP requests.}]
lower_jvp(target, wrt, tangent):
    target_binding = binding_of(target)
    wrt_binding = binding_of(wrt)
    seed = tangent or default_seed(wrt_binding, value=1)
    local_bindings = local_contexts[target_binding.defid]
    return differentiate_with_seed(target_binding.expr,
                                   wrt_binding,
                                   seed,
                                   local_bindings)

lower_vjp(target, wrt, cotangent):
    target_binding = binding_of(target)
    wrt_binding = binding_of(wrt)
    cot = cotangent or default_seed(target_binding, value=1)
    cot_defid, cot_ident = make_seed_placeholder(cot)

    objective = sum_over_target_axes(target * cot_ident)
    if target is scalar and wrt is tensor:
        bars = build_shared_scalar_target_bars(target, [wrt], cot)
        return simplify_and_hoist(bars[wrt.defid])

    if wrt is scalar:
        return substitute(cot_defid -> cot,
                          differentiate(objective, wrt, seed=1))

    wrt_axes = fresh axes for wrt shape
    basis = __basis_tensor(wrt.shape, wrt_axes)
    component = differentiate(objective, wrt, seed=basis)
    component = substitute(cot_defid -> cot, component)
    return EinsteinIR(indices=wrt_axes, value=simplify(component))
\end{lstlisting}

The scalar-target VJP branch is important for machine-learning examples: a loss
is often scalar while parameters are tensors. The lowerer factors the reverse
accumulation so shared intermediate bars can be reused instead of regenerating
the same derivative expression for every parameter index.

The JVP and VJP lowerers expose two complementary views of the same local
Jacobian. JVP seeds a parameter-direction tangent and asks how the target
changes; VJP seeds a target-direction cotangent and asks how that cotangent
pulls back to the parameter. The code keeps both paths in compile-time IR so
the generated derivative can still be simplified, hoisted, typed, and lowered
like any other expression.

The scalar-target optimization is more than a performance tweak. When a loss is
scalar and a parameter is a tensor, computing one derivative expression per
parameter element would duplicate the same upstream work. Building shared bars
for the scalar target lets the compiler reuse intermediate adjoints and then
project the requested parameter shape. This preserves the local source
interface while avoiding a naive dense-Jacobian expansion for common training
patterns.

The VJP construction uses a scalar objective because reverse-style propagation
is naturally a pullback from a cotangent. Multiplying the target by a cotangent
and summing over target axes turns a shaped target into a scalar whose
derivative with respect to the parameter has the parameter shape. This is the
compile-time IR equivalent of asking how a weighted observation of the target
changes with the parameter. The generated objective is ordinary source-like IR,
so simplification and hoisting can remove seed placeholders and common
subexpressions before backend execution.

The JVP construction is dual. A tangent seed with the parameter shape selects a
direction in parameter space, and differentiation pushes that direction forward
to the target. The lowerer keeps this as an explicit generated expression rather
than as a black-box runtime request, which means projected derivatives can still
benefit from tensor lowering and recurrence storage analysis. The two paths are
therefore not competing implementations; they are two projections of the same
Jacobian chosen according to how the source program indexes the lazy value.

\textit{Concrete example.} If \texttt{loss} is scalar and \texttt{W} is a
matrix, \texttt{@loss/@W} enters the scalar-target VJP branch. The lowerer
builds shared bars for the scalar objective and returns a matrix-shaped
gradient with the same axes as \texttt{W}. If instead \texttt{h[t]} is a vector
state and the program asks how a single parameter direction changes the whole
state, a basis tangent for that parameter selects a JVP column. In both cases
the seed is an ordinary generated tensor expression that subsequent range,
shape, and lowering passes can inspect.

\subsection{Lazy Jacobian Access Algorithm}

\begin{lstlisting}[style=pythoncode,caption={Projected lazy-Jacobian lowering.}]
rewrite_access(lazy[target, wrt], indices):
    target_shape = shape(target)
    wrt_shape = shape(wrt)
    out_rank = len(target_shape)
    wrt_rank = len(wrt_shape)

    out_indices = indices[:out_rank]
    wrt_indices = indices[out_rank:]

    if all output axes are fixed and no wrt axes are fixed:
        cotangent = basis_tensor(target_shape, out_indices)
        return lower_vjp(target, wrt, cotangent)

    if all wrt axes are fixed and mode(target, wrt) == "jvp":
        tangent = basis_tensor(wrt_shape, wrt_indices)
        column = lower_jvp(target, wrt, tangent)
        return column[out_indices]

    if all output axes are fixed and some wrt axes are fixed:
        row = lower_vjp(target, wrt, basis_tensor(target_shape, out_indices))
        return row[wrt_indices]

    return materialize_lazy_jacobian_as_einstein(target, wrt)
\end{lstlisting}

The materialization fallback still generates ordinary indexed IR. If JVP mode
is chosen, the generated Jacobian iterates over parameter axes, computes one
JVP column with a basis tangent, and indexes the target axes. If VJP mode is
chosen, it iterates over output axes, computes one VJP row with a basis
cotangent, and indexes the parameter axes.

The projection algorithm is where lazy Jacobians become operational. The
indices of a Jacobian access are split into target axes and parameter axes. If
the target side is fully fixed, the request is a row query and VJP is natural.
If the parameter side is fully fixed, the request is a column query and JVP is
natural. If both sides are partially fixed, the lowerer chooses the cheapest
known projection and indexes the result. Only the ambiguous case requires full
materialization.

This delayed decision is the main benefit of representing \texttt{@y/@x} as a
value. Source code can pass, index, or bind a derivative request without
committing to a dense matrix. The compiler observes the use site and then emits
ordinary IR for the demanded projection. That is why the lazy value still has a
shape: it behaves like a shaped tensor to the source language, even when the
backend computes only a row, a column, or a single entry.

The access splitter is deliberately syntactic after shape analysis. Once the
lazy Jacobian's shape is known, the first \texttt{out\_rank} indices always
refer to target axes and the remaining indices always refer to parameter axes.
That convention avoids ambiguity when target and parameter have the same rank or
similar index names. It also gives diagnostics a simple story: an out-of-bounds
or under-specified access can be explained in terms of the concatenated
Jacobian shape rather than in terms of a hidden matrix layout.

The fallback to materialization is still structured. The compiler does not call
into a generic dense-Jacobian runtime by default; it generates indexed IR that
iterates over the relevant axes and fills the lazy value. That fallback is more
expensive, but it remains analyzable by the normal compiler pipeline. This is an
important safety property: even when the projection optimizer cannot reduce the
request to a row or column, the program stays in the language's own IR rather
than escaping into an opaque autodiff subsystem.

\textit{Concrete example.} If \texttt{h} has shape \texttt{[T,H]} and
\texttt{theta} has shape \texttt{[P]}, then \texttt{J = @h/@theta} has shape
\texttt{[T,H,P]}. Accessing \texttt{J[t,j,:]} fixes all output axes and asks
for one row of the Jacobian, so lowering uses a VJP with a basis cotangent for
\texttt{h[t,j]}. Accessing \texttt{J[:,:,p]} fixes the parameter axis and asks
for one column, so lowering uses a JVP with a basis tangent for
\texttt{theta[p]}. Accessing \texttt{J[t,j,p]} can compute the row and then
index one parameter entry.

\section{Lazy Jacobians}

\pathcode{src/einlang/runtime/autodiff_requests.py} defines
\texttt{LazyJacobianValue}. A lazy Jacobian knows the target DefId, the
with-respect-to DefId, and the runtime autodiff engine. Its shape is the target
shape concatenated with the parameter shape. It chooses JVP mode when the
parameter dimension is no larger than the output dimension and VJP mode
otherwise. Indexing a single entry, row, or column can call a JVP or VJP without
materializing the full Jacobian.

This behavior is tested directly in \pathcode{tests/unit/test_lazy_jacobian.py}:
entry, row, and column indexing return correct results while leaving the dense
matrix unmaterialized.

\section{Runtime Autodiff Boundary}

The primary executed path is compile-time lowering: source \texttt{@} syntax is
removed before final lowered tensor execution. The runtime autodiff engine in
\pathcode{src/einlang/runtime/autodiff_requests.py} remains a thin internal
handler for autodiff builtins and for the standalone \texttt{LazyJacobianValue}
tests. It reads compiler-produced binding maps and local contexts when it needs
to evaluate a JVP or VJP request against the current backend environment.

\begin{lstlisting}[style=pythoncode,caption={Runtime lazy-value projection.}]
lazy_jacobian_getitem(item):
    split item into output key and wrt key
    if all output axes and all wrt axes are integers:
        if mode() == "jvp":
            return column(wrt_key)[output_key]
        else:
            return row(output_key)[wrt_key]
    if all output axes are integers:
        return row(output_key)[wrt_key]
    if all wrt axes are integers:
        return column(wrt_key)[output_key]
    return materialize()[item]
\end{lstlisting}

This runtime object is intentionally not the semantic center of autodiff. The
compiler path is responsible for eliminating source-level request nodes, and
\pathcode{src/einlang/passes/autodiff_leak_check.py} rejects any
\texttt{DifferentialIR}, \texttt{JvpIR}, \texttt{VjpIR}, or
\texttt{LazyJacobianIR} that survives past the allowed phase boundary.

The runtime projection code mirrors the compile-time lowering logic, but it is
used only when a lazy object reaches the runtime boundary in controlled tests or
internal builtin paths. It first separates output indices from parameter
indices, then chooses entry, row, column, or materialization behavior. The
important invariant is the same as in the compiler: indexing a Jacobian should
not force a full dense value unless the access pattern demands it.

Keeping this object small prevents two semantic centers from diverging. The
compiler owns source-level derivative meaning, while the runtime object
provides a practical lazy-value interface for projection behavior. The leak
check enforces the handoff. If a source \texttt{@} node survives too late, that
is a compiler bug rather than a backend feature to interpret.

The runtime lazy value exists mainly to make projection behavior concrete and
testable. Its mode choice is a heuristic over shape sizes: use JVP when
parameter-side directions are cheaper, and VJP when output-side cotangents are
cheaper. Because the compile-time lowering path owns ordinary source programs,
this heuristic does not define the language semantics. It defines how the
runtime object serves entry, row, column, and full requests when such an object
is intentionally constructed.

This boundary keeps future implementations flexible. A more advanced backend
could replace the runtime projection helper with compiled Jacobian kernels, or
the compiler could recognize more projection patterns statically. As long as the
phase invariant holds, those changes do not alter what \texttt{@y/@x} means in
source. The current implementation therefore uses the runtime object as a
practical bridge and a test target, while the thesis treats compile-time request
lowering as the semantic center.

\textit{Concrete example.} In the lazy-Jacobian unit tests, a runtime value for
\texttt{@y/@x} may be asked for \texttt{J[0,1]}, \texttt{J[0,:]}, or
\texttt{J[:,1]}. The first case computes a single entry through the cheaper of
row or column projection; the second computes a row; the third computes a
column. Only a slice that leaves both target and parameter axes unfixed calls
\texttt{materialize()}. This mirrors the compile-time rule and makes the
expected projection behavior executable in isolation.

\begin{table}[htbp]
\centering
\caption{Autodiff responsibilities by phase.}
\label{tab:autodiff-phases}
\renewcommand{\arraystretch}{1.08}
\begin{tabular}{@{}L{0.26\textwidth}L{0.62\textwidth}@{}}
\toprule
Phase & Responsibility \\
\midrule
High-level pass & Find source \texttt{@} forms, collect dependency facts, synthesize derivative IR, and preserve indexed structure. \\
Request lowering & Turn JVP, VJP, and lazy-Jacobian requests into ordinary IR; rewrite projected accesses when possible. \\
Leak check & Reject any source-level autodiff artifacts that survive past the allowed phase boundary. \\
Runtime engine & Support internal autodiff builtins and standalone lazy-value projection tests when such requests reach runtime. \\
Lazy value & Expose Jacobian shape, entry, row, column, and materialization behavior for runtime lazy objects. \\
\bottomrule
\end{tabular}
\end{table}

\section{Differential Rules and Builtins}

The source grammar supports \texttt{@fn} differential rules. The compiler merges
these into ordinary function definitions before name resolution so each function
has one DefId. Autodiff builtin requests are collected by scanning builtin calls
whose DefIds identify autodiff intrinsics. This lets derivative-related helpers
participate in ordinary IR and backend execution without requiring a separate
side language.

\section{Current Boundary}

Autodiff is broad enough for scalar arithmetic, elementwise tensor math,
reductions, indexed contractions, many standard-library functions, fitting
examples, and MNIST-style training examples. Unsupported derivative paths should
produce an error or conservative behavior rather than a silent finite-difference
approximation. The key implementation invariant is that source-level
\texttt{@} artifacts do not leak past the autodiff phase.
